\documentclass[11pt]{article}
\usepackage[margin=1in]{geometry}
\usepackage{amsmath,amssymb}
\usepackage{physics}
\usepackage{parskip}
\usepackage[hidelinks]{hyperref}

\title{Physics Derivations for Fiber Raman Suppression}
\author{}
\date{}

\newcommand{\uhat}{\hat{u}}
\newcommand{\utilde}{\tilde{u}}

\begin{document}
\maketitle

\section*{Scope}

Derivations unique to this project: the adjoint sensitivity equation for the
GMMNLSE, its discretization, the chain rule through the input phase mask, the
photon-number conservation law used for verification, and the SPM-corrected
time-window formula. The underlying GNLSE (Agrawal Eq.~2.3.36) and Blow--Wood
Raman response are taken as given.

\section{Adjoint sensitivity equation}

We want $\partial J / \partial \varphi(\omega)$ where $\varphi$ is the input
spectral phase and
\[
  J \;=\; \frac{E_\text{band}(L)}{E_\text{total}(L)}
      \;=\; \frac{\int_\mathcal{B} |\uhat(L,\omega)|^2 \,d\omega}
                  {\int |\uhat(L,\omega)|^2 \,d\omega},
\]
with $\uhat$ satisfying the GNLSE
\[
  \pdv{\uhat}{z} \;=\; F(\uhat,z)
  \;\equiv\; iD(\omega)\,\uhat
   + i\gamma\,\frac{\omega}{\omega_0}\,
     \mathcal{F}\!\left[(1-f_R)|u|^2 u + f_R(h_R \ast |u|^2)u\right],
\]
and initial condition $\uhat(0,\omega) = \uhat_0(\omega)\,e^{i\varphi(\omega)}$.

\subsection{Derivation via Lagrange multipliers}

Treat the ODE constraint as a Lagrangian equality and introduce an adjoint field
$\lambda(z,\omega)$:
\[
  \mathcal{L}
  \;=\; J[\uhat(L)]
  \;+\; \int_0^L \Re \left\langle \lambda,\,\pdv{\uhat}{z} - F(\uhat,z)\right\rangle dz.
\]
Here $\langle a,b\rangle = \int a^*(\omega) b(\omega)\,d\omega$. Integrate the
first term of the integral by parts in $z$:
\[
  \int_0^L \Re \left\langle \lambda, \pdv{\uhat}{z}\right\rangle dz
  \;=\; \Re \langle \lambda(L),\uhat(L)\rangle - \Re \langle \lambda(0),\uhat(0)\rangle
        - \int_0^L \Re \left\langle \pdv{\lambda}{z},\uhat\right\rangle dz.
\]
Requiring stationarity of $\mathcal{L}$ with respect to variations
$\delta\uhat(z)$ for $z\in(0,L)$ gives
\begin{equation}
  -\dv{\lambda}{z} \;=\; \left[\pdv{F}{\uhat}\right]^\dagger \lambda,
  \label{eq:adjoint_ode}
\end{equation}
stationarity at the endpoint $z=L$ gives the terminal condition
\begin{equation}
  \lambda(L,\omega) \;=\; \pdv{J}{\uhat^*(L,\omega)},
  \label{eq:terminal}
\end{equation}
and stationarity at $z=0$ produces the sensitivity with respect to the initial
condition,
\begin{equation}
  \pdv{J}{\uhat^*(0,\omega)} \;=\; \lambda(0,\omega).
  \label{eq:sensitivity_u0}
\end{equation}
Eq.~(\ref{eq:adjoint_ode}) is integrated backward from $L$ to $0$ with terminal
data Eq.~(\ref{eq:terminal}).

\subsection{Terminal condition for our cost}

Differentiate $J = E_\text{band}/E_\text{total}$ with respect to $\uhat^*(L,\omega)$
using $\partial |\uhat|^2 / \partial \uhat^* = \uhat$:
\begin{align*}
  \pdv{J}{\uhat^*(L,\omega)}
  &= \frac{\mathbb{1}_\mathcal{B}(\omega)\,\uhat(L,\omega)\,E_\text{total}
            - \uhat(L,\omega)\,E_\text{band}}
           {E_\text{total}^2} \\[0.3em]
  &= \frac{\uhat(L,\omega)\,[\mathbb{1}_\mathcal{B}(\omega) - J]}{E_\text{total}},
\end{align*}
where $\mathbb{1}_\mathcal{B}$ is the indicator function of the Raman band.
Therefore
\begin{equation}
  \lambda(L,\omega)
  \;=\; \frac{\uhat(L,\omega)\,\bigl[\mathbb{1}_\mathcal{B}(\omega) - J\bigr]}
              {E_\text{total}}.
  \label{eq:terminal_explicit}
\end{equation}
Implemented in \texttt{scripts/common.jl} L269 as
\texttt{dJ = u\textomega f .* (band\_mask .- J) ./ E\_total}.

\subsection{Adjoint operator in the interaction picture}

The forward ODE is integrated for the interaction-picture field
$\utilde(z,\omega) = e^{-iD(\omega)z}\,\uhat(z,\omega)$, so the adjoint is
integrated for the corresponding interaction-picture multiplier
$\tilde{\lambda}(z,\omega) = e^{-iD(\omega)z}\,\lambda(z,\omega)$. The terminal
condition becomes $\tilde{\lambda}(L,\omega) = e^{-iD(\omega)L}\lambda(L,\omega)$.
This transformation is applied in
\texttt{src/simulation/sensitivity\_disp\_mmf.jl} L296.

The operator $[\partial F/\partial \uhat]^\dagger$ has four contributions
because the nonlinearity $|u|^2 u$ has two derivatives (with respect to both $u$
and $u^*$ via $|u|^2 = u u^*$) and the Raman convolution has a forward and
conjugate path. Let $v = \Re u$, $w = \Im u$ be the forward field in time, and
let $\delta_K^{(1)}$, $\delta_K^{(2)}$ be the two Kerr operators
\begin{align}
  \delta_K^{(1)}[t,i,j] &= \sum_{k\ell} \gamma_{\ell k i j}\,(v_k v_\ell + w_k w_\ell)
      &&(\text{from } |u|^2 = u u^*) \label{eq:delta1}\\
  \delta_K^{(2)}[t,i,j] &= \sum_{k\ell} \gamma_{\ell k i j}\,(v_k v_\ell - w_k w_\ell + 2i v_k w_\ell)
      &&(\text{from } u^2) \label{eq:delta2}
\end{align}
and let $\delta_R^{(1)} = h_R \ast \delta_K^{(1)}$ (the Raman convolution of the
$|u|^2$ operator). The Kerr+Raman combined adjoint operator needs the factor
of 2 from $\partial(|u|^2 u)/\partial u^* = u^2$, giving
\begin{equation}
  \delta_{KR} \;=\; 2(1-f_R)\,\delta_K^{(1)} \;+\; \delta_R^{(1)}.
  \label{eq:deltaKR}
\end{equation}
The adjoint RHS in the interaction picture is then
\begin{equation}
  \dv{\tilde{\lambda}}{z}
  \;=\; i\,e^{-iD(\omega)z}\,\mathcal{F}^{-1}\!\Bigl[\,
          \lambda_t \cdot \delta_{KR}
          \;+\; (1-f_R)\,\lambda^*_t \cdot \delta_K^{(2)}
          \;+\; R_\text{fwd}(\lambda)
          \;+\; R_\text{conj}(\lambda)
        \,\Bigr],
  \label{eq:adjoint_rhs}
\end{equation}
where $R_\text{fwd}$ and $R_\text{conj}$ are the two Raman-convolution paths
that arise from the self-steepening-scaled Raman term acting on $\lambda$ and
$\lambda^*$ respectively. These paths involve a convolution with $h_R^*$ or
$h_R$ plus a $\gamma$-tensor contraction with the forward field at $z$.

Implemented in \texttt{sensitivity\_disp\_mmf.jl} L32--79 as four helper
functions summed on L78:
\begin{itemize}
\item L72 \texttt{calc\_$\lambda$\_$\partial$fKR1c$\partial$uc!} computes the
  $\lambda_t \cdot \delta_{KR}$ term
\item L73 \texttt{calc\_$\lambda$c\_$\partial$fK$\partial$uc!} computes the
  $\lambda^*_t \cdot \delta_K^{(2)}$ term
\item L74 \texttt{calc\_$\lambda$\_$\partial$fR2c$\partial$uc!} computes $R_\text{fwd}$
\item L75 \texttt{calc\_$\lambda$c\_$\partial$fR$\partial$uc!} computes $R_\text{conj}$
\end{itemize}

\subsection{Chain rule through the input phase mask}

The input field is $\uhat(0,\omega) = \uhat_0(\omega)\,e^{i\varphi(\omega)}$.
We want $\partial J / \partial \varphi(\omega)$, but we already have
$\partial J / \partial \uhat^*(0,\omega) = \lambda(0,\omega)$ from
Eq.~(\ref{eq:sensitivity_u0}). Since $J$ is real-valued,
\[
  dJ \;=\; 2\,\Re\!\left\{\int \lambda^*(0,\omega)\,d\uhat(0,\omega)\,d\omega\right\}.
\]
For a phase perturbation $d\varphi$,
$d\uhat(0,\omega) = \uhat_0(\omega)\,i e^{i\varphi(\omega)}\,d\varphi(\omega)
= i\,\uhat_\text{shaped}(\omega)\,d\varphi(\omega)$, so
\begin{equation}
  \pdv{J}{\varphi(\omega)}
  \;=\; 2\,\Re\!\bigl[\,\lambda^*(0,\omega)\cdot i\,\uhat_\text{shaped}(\omega)\,\bigr].
  \label{eq:phase_gradient}
\end{equation}
Implemented in \texttt{scripts/raman\_optimization.jl} L91 as
\texttt{$\partial$J\_$\partial\varphi$ = 2.0 .* real.(conj.($\lambda$0) .* (1im .* u$\omega$0\_shaped))}.

\subsection{Verification via Taylor remainder}

If $\nabla J$ is the true gradient, Taylor's theorem gives
\begin{equation}
  R(\epsilon) \;\equiv\;
  \bigl|J(\varphi + \epsilon\,\delta\varphi) - J(\varphi) - \epsilon\,\langle \nabla J,\delta\varphi\rangle\bigr|
  \;=\; \tfrac{1}{2}\epsilon^2\,|\langle H\delta\varphi,\delta\varphi\rangle| + O(\epsilon^3).
  \label{eq:taylor}
\end{equation}
So $\log R$ versus $\log \epsilon$ has slope $2$. A linear bias in the gradient
$\nabla J_\text{claimed} = \nabla J + b$ produces an extra term
$\epsilon\,|\langle b,\delta\varphi\rangle|$ in $R$ which dominates at small
$\epsilon$ and drops the slope to $1$. Measured slopes at $N_t = 2^{14}$,
$L = 0.1$ m, $\epsilon \in \{10^0,10^{-1},10^{-2},10^{-3}\}$: $2.01$, $2.07$,
$2.09$.

\section{Photon number conservation}

With self-steepening and Raman, \emph{classical energy} $\int |\uhat(\omega)|^2 d\omega$ is
not conserved because every Raman event transfers energy $\hbar(\omega_1-\omega_2)$
into a phonon. The conserved quantity is the photon number. Each photon at
$\omega$ carries energy $\hbar\omega$, so the spectral photon density is
$|\uhat(\omega)|^2/(\hbar\omega)$ and
\begin{equation}
  N_\text{ph} \;=\; \int \frac{|\uhat(\omega)|^2}{\hbar\,\omega}\,d\omega.
  \label{eq:photon_number}
\end{equation}
A Raman event at $\omega_1 \to \omega_2$ removes energy $\hbar\omega_1$ at
$\omega_1$ (hence $\hbar\omega_1/\hbar\omega_1 = 1$ photon) and deposits energy
$\hbar\omega_2$ at $\omega_2$ (hence $\hbar\omega_2/\hbar\omega_2 = 1$ photon),
leaving $N_\text{ph}$ invariant. The $\hbar$ cancels in the ratio
$N_\text{ph}(L)/N_\text{ph}(0)$ so we drop it in code. Implemented in
\texttt{scripts/verification.jl} L189--199 as
\texttt{sum(abs2.(u$\omega$) ./ abs.(omega\_s)) * $\Delta$t}.

Measured drift across five production runs after the SPM time-window fix:
all $< 1\%$. Prior to the fix, high-power runs showed up to $49\%$ drift,
which pointed directly to the dimensional error in the next section.

\section{Time-window sizing with SPM broadening}

The time window $W$ must contain the pulse throughout propagation. Two effects
broaden the pulse in time.

\paragraph{Dispersive walk-off.} Spectral components separated by the Raman
shift $\Delta\omega_R = 2\pi \cdot 13\,\mathrm{THz}$ travel at different group
velocities, so they separate by
\begin{equation}
  \Delta t_\text{walk} \;=\; |\beta_2|\,L\,\Delta\omega_R.
  \label{eq:walkoff}
\end{equation}

\paragraph{SPM-induced spectral broadening transmitted through GVD.} The
nonlinear phase accumulated over the fiber is
\begin{equation}
  \phi_\text{NL} \;=\; \gamma\,P_\text{peak}\,L \qquad [\text{radians, dimensionless}].
  \label{eq:phi_NL}
\end{equation}
For a sech$^2$ pulse, the SPM-broadened angular bandwidth is
\begin{equation}
  \Delta\omega_\text{SPM} \;=\; 0.86\,\frac{\phi_\text{NL}}{T_0}
  \qquad [\text{rad/s}],
  \label{eq:domega_spm}
\end{equation}
where $T_0 = \text{FWHM}/1.763$ (Agrawal Ch.~4). Propagating this broadened
spectrum through GVD gives an additional temporal spread
\begin{equation}
  \Delta t_\text{SPM} \;=\; |\beta_2|\,L\,\Delta\omega_\text{SPM}.
  \label{eq:dt_spm}
\end{equation}

\paragraph{Dimensional point.} $\phi_\text{NL}$ in Eq.~(\ref{eq:phi_NL}) is
\emph{dimensionless}. Using it directly as an angular frequency is dimensionally
wrong: the $1/T_0$ factor in Eq.~(\ref{eq:domega_spm}) converts radians to
rad/s. An earlier version of the code was missing this $1/T_0$ and therefore
undersized the time window by a factor of order $T_0 \cdot (\text{scale})$,
which at 185\,fs pulses and typical operating powers was 5--7$\times$. The
photon conservation test in the previous section caught this (49\% drift at
high power) and the correction restored conservation to $<1\%$.

Implemented in \texttt{scripts/common.jl} L191--215: L208 computes
$\phi_\text{NL}$, L209 divides by $T_0$ to get Eq.~(\ref{eq:domega_spm}), L210
applies Eq.~(\ref{eq:dt_spm}), and L213--214 sums walk-off plus SPM with a
safety factor of 2.

\end{document}
